\chapter{Examples with prescribed targets}\label{tx:ch:examples}
The classification allows both finite proof degrees and infinity. An
example should make clear which kind of hardness it illustrates. We
already have a finite T2 family from a Boolean first-order sentence:
the descending-predecessor example of
Chapter~\ref{tx:ch:running-order}, with its full lifting proof. Here we
begin with a consistent family whose target is derivable on one residue
class and underivable on the others. We then obtain finite hard proofs
of a nonconstant target. Finally we describe the four first-order
refutation families used by the undecidability argument.
Recall the test in Theorem~\ref{tx:thm:classification}: T1 is eventually
bounded; T2 has no recurring finite degree; T3 is not eventually bounded and has a finite
nonempty set of recurring finite degrees; T4 has infinitely many. For a request,
$h_S$ is its least-degree profile; $d_S(\Gamma)$ denotes the special
least degree of a refutation, whose target is $1$.

\section{A consistent family with an exact residue split}
Fix characteristic $p$ and, for $n\ge1$, take
\[
                  \Gamma_n=\{x_i-1:i<n\},\qquad
                  G_n=\sum_{i<n}x_i.
\]
The axioms have the solution $x_i=1$ for every $i$, so they are
consistent at every size. Both the axioms and the invariant target
have stable coefficient descriptions.

\begin{proposition}\label{tx:prop:consistent-T3}
For both NS and PC, the target profile is exactly
\[
 h_S(n)=\begin{cases}1,&p\mid n,\\\infty,&p\nmid n.\end{cases}
\]
In particular, it is T3. With the same axioms but the indexed request
$G_{i,n}=x_i$ for every $i<n$, the profile is constantly infinity and
has type T2.
\end{proposition}
\begin{proof}
Summing the axioms gives $G_n-n$. If $p\mid n$, this is $G_n$, so it
has an NS certificate, and hence a PC proof, of degree one. The target
is nonzero of degree one, proving equality. If $p\nmid n$, the root
with all coordinates one evaluates $G_n$ to the nonzero field element
$n$. Soundness excludes any derivation. The same root evaluates each
indexed target $x_i$ to one, proving the final assertion.
\end{proof}

This example also distinguishes an invariant sum from its individual
summands. The statement that the sum is a consequence does not say
that each summand is a consequence. The all-target profile is a maximum
of separate proof degrees, not the degree of a proof of their sum.

The easy residue class here is zero modulo $p$; its degree is exactly
one. The other classes have infinite degrees. This is a genuine
consistent-axiom example of the classification, but it does not yet
provide finite proofs requiring degree $\Omega(n)$. The next
construction addresses that different question.

\section{A fresh variable as a nonconstant target}\label{tx:sec:target-examples}
The degree of a target need not reflect the degree needed to derive it.
A polynomial of degree one may require arbitrarily high proof degree.
We first give an exact construction, then apply it to a uniform family
whose hard values are finite.

A simple invariant target illustrates the distinction between finite
hardness and failure of derivability. For $G_n=\sum_{i<n}x_i$, the input
$\{1\}$ gives degree exactly one in both systems for $n\ge1$. With no
inputs the same target is underivable and has degree $\infty$. Both
fit the uniform definition, but neither supplies a finite hard example.
The next construction does.

Here \emph{fresh} means a new indeterminate in the polynomial ring,
independent of every old variable and absent from every old axiom.
Any coefficient polynomial has a unique expansion
$a(x,y)=\sum_j a_j(x)y^j$; the notation $[y^1]a$ selects $a_1(x)$.
This linear operation is different from substituting a value for $y$.
The proof uses substitution for PC and coefficient extraction for NS,
because they control the two degree conventions differently.

\begin{proposition}[Adjoining a fresh target]\label{tx:prop:fresh-target}
Let $K$ be any field, let $\Gamma\subseteq K[x_v:v\in V]$ be any axiom
set, and adjoin a fresh variable $y$. Each certificate uses only
finitely many axioms. Let $\Gamma^\uparrow$ be the same axioms viewed
in $K[x_v:v\in V,y]$, with no additional axiom involving $y$. Then
\begin{equation}\label{tx:eq:fresh-identities}
 \begin{aligned}
 h_{\PC}(\Gamma^\uparrow\vdash y)
       &=\max\{1,d_{\PC}(\Gamma)\},\\
 h_{\NS}(\Gamma^\uparrow\vdash y)
       &=d_{\NS}(\Gamma)+1.
 \end{aligned}
\end{equation}
These equalities include degree zero and infinity, with
$\infty+1=\infty$.
\end{proposition}
\begin{proof}
For PC, embed a refutation of $\Gamma$ in the larger ring and multiply
its last line $1$ by $y$. If the refutation had degree $d$, the resulting
proof of $y$ has degree $\max\{1,d\}$: only the new last line is added,
and its degree is one.

Conversely, substitute $y=1$ in every line of a PC proof of $y$.
The old axioms are unchanged and linear combinations commute with the
substitution. Multiplication by an old variable remains a multiplication
step; multiplication by $y$ becomes a repeated line, which the linear
combination rule admits. No degree increases, and the last line becomes
$1$. Every proof of $y$ also has degree at least one. These two lower
bounds prove the first equality, including the case when no refutation
exists.

For NS, multiplying the coefficients of an identity
$1=\sum_fa_ff$ by $y$ gives an identity for $y$ with degree one larger.
For the reverse inequality, suppose
\[
                    y=\sum_f a_f(x,y)f(x)
\]
has total-product degree at most $e$. Take the coefficient of $y^1$.
Since every $f$ is independent of $y$, this gives the exact identity
\[
                    1=\sum_f [y^1]a_f(x,y)\,f(x).
\]
Every nonzero product on the right has degree at most $e-1$: it is the
coefficient of $y^1$ in a product of total degree at most $e$. Zero
products may be discarded. Hence $d_{\NS}(\Gamma)\le e-1$.
In particular a finite proof of $y$ exists exactly when a finite
refutation exists. If the original refutation degree is zero, both
new degrees are one, as asserted.
\end{proof}

\begin{example}[Why the two exact identities differ]
Take $\Gamma=\{x,1-x\}$. Both refutation degrees are one. PC derives
$1=x+(1-x)$ and then multiplies that last constant by $y$; every line
still has degree at most one. An NS identity for the fresh target is
\[
                             y=yx+y(1-x).
\]
Its products have degree two. No degree-one NS identity can produce
$y$, since its multipliers of these degree-one axioms would have to
be constants. Thus the new PC degree is one and the NS degree is two.
PC can finish the refutation before multiplying its last line; an NS
identity carries the extra variable in its coefficient polynomials.
\end{example}

The absence of a $y$-axiom is part of the construction; in particular,
we do not add $y^2-y=0$. For a uniform template, give the new variable
empty support and keep all old axiom indices and coefficients. Set
all coefficients of axiom monomials involving $y$ to zero. The result
is again a uniform template with the same bounds $k,s,D,n_0$, and
$G_n=y$ is a fully invariant target of degree one at every rank.
Thus \eqref{tx:eq:fresh-identities} holds rank by rank.

This operation preserves each of the four types, separately for NS
and PC. The NS recurring finite degrees are shifted by one. For PC,
the map $d\mapsto\max\{1,d\}$ merges only degrees zero and one;
all larger finite degrees and infinity are unchanged. Eventual
boundedness, absence or unboundedness of recurring finite levels, and
the mixed periodic case are consequently preserved.

The notation $v_p(n)$ is the exponent of $p$ in the factorisation of
$n$: writing $n=p^a u$ with $p\nmid u$ gives $v_p(n)=a$.
For instance, $v_3(12)=1$ and its associated power is $q=3$.
A window such as $[q,q+4]$ bounds a finite proof degree; the condition
$n\ge2q$ must still be checked before using its lower endpoint.

\begin{corollary}[A finite hard nonconstant consequence family]\label{tx:cor:fresh-peaks}
Let $K$ have prime characteristic $p$, and let $\widehat B_n$ be the
normalised weak-bijection system defined in
\eqref{tx:eq:bhat-witness}, with all of its function and coverage witnesses
retained. Adjoin a fresh variable $y$ as above, and request the target
$G_n=y$. Both target-degree profiles are of type T4.

More precisely, write $n=p^a u$, where $p\nmid u$, and put $q=p^a$.
Whenever $n\ge2q$,
\[
 \begin{gathered}
 q\le h_{\PC}(n)\le q+4,\qquad
 q+1\le h_{\NS}(n)\le q+5,\\
                         h_{\PC}(n)\le h_{\NS}(n).
 \end{gathered}
\]
In particular, at every size $n_a=(p+1)p^a$,
\begin{equation}\label{tx:eq:fresh-peaks}
 \frac{n_a}{p+1}\le h_{\PC}(n_a)\le h_{\NS}(n_a)
                              \le\frac{n_a}{p+1}+5<\infty.
\end{equation}
\end{corollary}
\begin{proof}
The refutation bounds \eqref{tx:eq:bhat-window}, proved below together with the rectangular obstruction in the next
chapter, give
$q\le d_{\PC}(\widehat B_n)\le d_{\NS}(\widehat B_n)\le q+4$.
Apply \eqref{tx:eq:fresh-identities}; since $q\ge1$, the maximum in the
PC identity changes nothing here. At $n_a=(p+1)p^a$ we have
$v_p(n_a)=a$ and $n_a\ge2p^a$, yielding \eqref{tx:eq:fresh-peaks}.

For each fixed $a$, infinitely many sizes
$n=p^a(p(j+1)+1)$ have the same valuation and satisfy the displayed
bounds. Some finite degree in each of the bounded windows must
therefore recur infinitely often. The windows tend to infinity with
$a$, proving T4 in each system. On the other hand, the old upper
bound gives $h_{\PC}(n)\le5$ and $h_{\NS}(n)\le6$ whenever
$p\nmid n$. The linear peaks thus coexist with infinitely many easy
sizes, rather than giving an eventual linear lower bound at all sizes.
\end{proof}

This example has finite hard proofs of a nonconstant target, but its
axioms are inconsistent: the difficulty is inherited from refuting
$\widehat B_n$. It does not establish a family of satisfiable axioms
with the same hard consequence profile. The small consistent example
in Chapter~\ref{tx:ch:consequences} illustrates ordinary-PC consequences in
a different setting; no growing lower bound is asserted for it.

\section{The refutation special case: four first-order examples}
\label{tx:sec:examples}

We now take the constant target $G_n=1$. This is the refutation special
case of the uniform-consequence theorem; the sharper invariant-target
budget applies. Write $\psi_1,\psi_2,\psi_3,\psi_4$ for the four sentences described in order below.
The four cases all occur among fixed first-order sentences with no finite
nonempty models. We give examples over a field of fixed prime characteristic
$p$. Each sentence is translated by the graph-and-sum convention: relation
and function-graph variables are Boolean, and every function has a totality
sum and pairwise functionality equations, as in
Chapter~\ref{tx:sec:systems}. These are ordinary polynomial inputs,
including the Boolean equations. The parameter $n\geq1$ is the
universe size. For the matrix notation in this section, we relabel the
universe as $\{1,\ldots,n\}$; this changes no equation or degree.

\paragraph{T1: a local contradiction.}
The sentence $\forall x\,(R(x)\wedge\neg R(x))$ contributes $r_i$ and
$1-r_i$. Their sum refutes the system in degree one, independently of $n$.
The contradiction is already visible at one element.

\paragraph{T2: an injection that misses a point.}
Let $\eta$ say that a unary function $g$ is injective and that a constant
$c$ is outside its image. Write $x_{ij}$ for $g(i)=j$ and $c_j$ for $c=j$.
In addition to the Boolean equations, its system $\Gamma_n^\eta$ consists of
\begin{equation}\label{tx:eq:G-witness}
\begin{gathered}
 \sum_jx_{ij}-1=0,\qquad
 x_{ij}x_{ik}=0\ (j\ne k),\qquad
 x_{ij}x_{hj}=0\ (i\ne h),\\
 \sum_jc_j-1=0,\qquad
 c_jc_k=0\ (j\ne k),\qquad x_{ij}c_j=0.
\end{gathered}
\end{equation}
A solution would inject $n$ elements into the $n-1$ elements other than $c$,
so every instance is unsatisfiable. For either proof system, every
degree-$d$ refutation satisfies
\[
 n<3pd\qquad(n\geq3).
\]
Here is how the internal matching bound gives this estimate. Choose the
largest $p$-power $q\leq n/3$ and put $r=n-q$. Set $c_n=1$, all other
$c_j=0$, and restrict the image of $g$ to $[r]$. After transposition, every
remaining axiom is an axiom of the bijection system $B_{r,n}$.
Here $B_{r,m}$ is the Boolean $r$-by-$m$ matrix system with every row and
column sum equal to one and both row and column collisions equal to zero.
If $d<q$, then $2d\leq r$ and $q\mid n-r$, contradicting the rectangle
obstruction of Lemma~\ref{tx:lem:rectangle-obstruction}. The substitutions do
not increase ordinary-PC or NS degree. Thus $d\geq q$, while maximality of
$q$ gives $n<3pq\leq3pd$. This proves T2 without an external pigeonhole
lower-bound theorem.

\paragraph{The injection certificate and the linear scale.}
The paper's matching upper bound is useful to inspect directly. Put
\[
 R_i=\sum_{j=1}^n x_{ij},\qquad C=\sum_{j=1}^n c_j,
                    \qquad P=\prod_{i=1}^n R_i.
\]
Every monomial in $CP$ has the form $c_j\prod_i x_{i,\lambda(i)}$.
If $\lambda$ is not injective, two factors give a column-collision
axiom. If it is injective, it is a permutation of $[n]$, and its
unique row $i$ with $\lambda(i)=j$ gives the omission axiom
$c_jx_{ij}$. In either case the monomial is an axiom multiple of
degree $n+1$. Summing these identities expresses $CP$ in the raw NS
span at that degree.

The second certificate is the exact telescoping identity
\[
 CP-1=(C-1)+C\sum_{i=1}^n(R_i-1)\prod_{h<i}R_h.
\]
Its summands use only constant and row totality axioms, with each
complete product of degree at most $n+1$. Subtracting this certificate
from the one for $CP$ derives $1$ at degree at most $n+1$.
For a two-element universe the telescoping list is
$(C-1)+C(R_1-1)+CR_1(R_2-1)$, with degrees $1,2,3$.
A term $c_1x_{11}x_{21}$ in $CP$ is a column-collision multiple;
$c_1x_{11}x_{22}$ is an omission-axiom multiple. These are the two
possibilities in the general expansion.

Thus, over any field and for every $n\ge1$,
\[
 d_{\PC}(\Gamma_n^\eta)\le d_{\NS}(\Gamma_n^\eta)\le n+1.
\]
The PC inequality uses the degree-preserving simulation of an NS
certificate by repeated variable multiplication. No Boolean reduction
is used. With the preceding $n<3pd$ lower bound for $n\ge3$, both
degrees are $\Theta(n)$ for each fixed positive characteristic.
This family proves sharpness of the general linear scale; it does
not supply an upper bound for every uniform family. This certificate
and the sharpness deduction are written mathematics, as in the paper;
no additional Lean endpoint is claimed for them.

\paragraph{T3: hardness on one residue class.}
Let $\chi_p$ assert
\[
 \forall x\qquad f^p(x)=x\quad\hbox{and}\quad f(x)\ne x.
\]
Its finite models are disjoint unions of directed $p$-cycles. In particular,
it has a $0/1$ graph assignment exactly when $p\mid n$. An explicit
normalisation uses the function equations for $f$, the equations $f_{ii}=0$,
and
\begin{equation}\label{tx:eq:cycle-path}
 f_{i_0i_1}f_{i_1i_2}\cdots f_{i_{p-1}i_p}=0
 \qquad(i_p\ne i_0).
\end{equation}
If $p\nmid n$, summing all length-$p$ walks gives an NS refutation of degree
at most $p$. Indeed, totality derives the sum of their monomials minus $n$
within degree $p$. Nonclosed walks vanish by \eqref{tx:eq:cycle-path}.
A closed walk with a repeated vertex either contains a functionality
collision or follows a shorter deterministic cycle; since $p$ is prime,
the latter is a forbidden fixed point. Each remaining cycle contributes its
monomial $p$ times, according to its starting vertex, and hence contributes
zero. The resulting nonzero scalar $n$ yields the refutation.

Now take $\psi_3=\eta\wedge\chi_p$ in disjoint nonlogical symbols.
For $p\mid n$, substitute a $p$-cycle assignment into the second block.
This preserves proof degree and shows that the conjunction has exactly the
degree of $\eta$, separately for NS and PC. For $p\nmid n$, the second
block refutes in degree at most $p$. Consequently this all-size
unsatisfiable family is bounded off the multiples of $p$ and has linear
degree on the multiples: it has type T3. No exact degree is needed on the
bounded residue classes.

\paragraph{T4: the valuation of the universe size controls the difficulty.}
The weak-bijection example is more revealing. Let unary functions $f,g$
be injective with disjoint images, and require their images together to
cover the universe. Use a unary witness $h$ for coverage:
\[
 \forall z\qquad f(h(z))=z\ \vee\ g(h(z))=z.
\]
Let $a_{ij},b_{ij},w_{ji}$ represent $f(i)=j$, $g(i)=j$, and $h(j)=i$.
The combined $2n$-by-$n$ matrix $(a;b)$ has row sums equal to one and both
row and column collisions equal to zero. The witness block consists of
\begin{equation}\label{tx:eq:bhat-witness}
 \sum_iw_{ji}-1=0,\qquad
 w_{ji}w_{ji'}=0\ (i\ne i'),\qquad
 w_{ji}(1-a_{ij})(1-b_{ij})=0,
\end{equation}
together with the Boolean equations. Call this complete system
$\widehat B_n$. Its witnesses are retained: it is not identified with the
simpler matrix system $B_{n,2n}$.

Put $q=p^{v_p(n)}$. The proved bounds give
\begin{equation}\label{tx:eq:bhat-window}
 q\leq d_{\mathrm{PC}}(\widehat B_n)
 \leq d_{\mathrm{NS}}(\widehat B_n)\leq q+4
 \qquad\hbox{whenever } n\geq2q.
\end{equation}
The upper bound holds at every positive size, without the displayed size
condition. We give both degree-preserving reduction and certificate
explicitly.

A substitution is used here as a map of proofs, not merely as a map
of solutions. An affine substitution replaces a variable by a linear
polynomial plus a constant, so it does not increase polynomial degree.
One must also check that every substituted input is derivable from the
new inputs within the permitted degree. The following identities supply
that check for the coverage and Boolean witness equations.

For the lower bound transpose the combined matrix and substitute
$w_{ji}=a_{ij}+b_{ij}$ into every witness variable. Its totality becomes
the corresponding totality of $B_{n,2n}$. The identities
\begin{align*}
(a+b)^2-(a+b)&=(a^2-a)+(b^2-b)+2ab,\\
(a+b)(1-a)(1-b)&=-(1-b)(a^2-a)-(1-a)(b^2-b)
\end{align*}
put the substituted Boolean and coverage equations in the raw NS spans
at degrees two and three respectively. A substituted witness collision
expands into four column collisions of the raw combined matrix.
Thus every substituted input has a derivation within its own degree.
Substitution is affine: a multiplication by an old variable becomes a
sum of permitted multiplications by new variables, with the preceding
line still of degree at most $d-1$. A degree-$d$ ordinary-PC refutation
would therefore refute $B_{n,2n}$ within the same degree. The rectangle
obstruction excludes this when $2d\le n$ and $d<q$, yielding the lower
bound in \eqref{tx:eq:bhat-window}.

For the upper bound fix a column $j$ and abbreviate
\[
 u_i=a_{ij}+b_{ij},\qquad T=\sum_i u_i,\qquad
 Z=\sum_i w_{ji}u_i.
\]
The coverage axiom and disjointness collision give the exact identity
\[
 Z-1=(\sum_iw_{ji}-1)
 -\sum_iw_{ji}(1-a_{ij})(1-b_{ij})
 +\sum_iw_{ji}a_{ij}b_{ij}.
\]
Consequently $Z-1$ lies in the degree-three NS span. If $x$ is one entry
of this column, then
\[
 Tx-x=(x^2-x)+\sum_{y\ne x}yx
\]
lies in the degree-two span, by Booleanity and column collisions.
Adding this for $a_{ij}$ and $b_{ij}$ gives $Tu_i-u_i$ in that span.
Finally
\[
 T-1=\sum_iw_{ji}(Tu_i-u_i)-(T-1)(Z-1)
\]
puts the missing column totality in the degree-four span. This retains
the coverage witnesses throughout the derivation.

For completeness, consider any matrix with row sums $R_i$ and column
sums $C_j$, with their differences from one in the NS span at degree
$e$, and with both collision families in degree two. For a fixed
$k$-element row set $S$, telescoping the product of its row sums gives
$\prod_{i\in S}R_i-1$ in degree at most $e+k$. Summing over $S$ gives
\[
 \sum_{|S|=k}\prod_{i\in S}R_i-\binom{\#\mathrm{rows}}k
\]
in that same span. Expand the products as partial functions from $k$
rows to columns. A noninjective partial function contains a column
collision, so its monomial is already in the degree-$k$ span.
The remaining monomials are exactly those of $k$-edge matchings.
The column calculation gives the same matching sum, since inversion
pairs partial injections with identical monomials. Subtracting derives
\[
 \binom{\#\mathrm{rows}}k-\binom{\#\mathrm{columns}}k
\]
within degree $e+k$. Applied to the $2n$-by-$n$ matrix with $e=4$, this
is the desired counting certificate.

Write $n=qu$ with $q=p^{v_p(n)}$ and $p\nmid u$. In characteristic
$p$, $(1+t)^n=(1+t^q)^u$. The coefficient of $t^q$ is therefore $u$;
for $(1+t)^{2n}$ it is $2u$. Thus
$\binom{2n}q-\binom nq=u\ne0$ in $K$. Rescaling the counting
certificate gives an NS refutation of degree at most $q+4$.

For each fixed $a$, the sizes $n=p^a(p(k+1)+1)$, $k\geq0$, have
$v_p(n)=a$ and satisfy \eqref{tx:eq:bhat-window}. Hence some finite degree in
$[p^a,p^a+4]$ recurs infinitely often. These windows are unbounded as $a$
varies, proving T4 for both proof systems. On the other hand, $p\nmid n$
implies degree at most five. There are explicit simultaneous linear peaks:
at $n_a=(p+1)p^a$,
\[
 \frac{n_a}{p+1}\leq d_{\mathrm{PC}}(\widehat B_{n_a})
 \leq d_{\mathrm{NS}}(\widehat B_{n_a})
 \leq\frac{n_a}{p+1}+4.
\]
For example, in characteristic two every odd $n$ has degree at most
five, whereas at $n=3\cdot2^a$ both degrees lie between $2^a$ and
$2^a+4$. The peaks are linear in their universe size even though the
profile returns to bounded values infinitely often.
Thus linear peaks coexist with infinitely many easy sizes. This is the
reason the T4 conclusion concerns peaks rather than an eventual linear
lower bound at every size.

\section{Exercises}
\begin{exercise}\label{tx:ex:sum-versus-individual}
For the consistent family of Proposition~\ref{tx:prop:consistent-T3},
compute the two profiles at sizes $1,2,3,4,5,6$ in characteristic three.
Compare the invariant sum request with the indexed request for every
individual variable. Which statements have finite proof degrees?
\end{exercise}
\noindent{\small\hyperref[tx:ans:sum-versus-individual]{Model answer (p.~\pageref*{tx:ans:sum-versus-individual}).}}\par\medskip

\begin{exercise}\label{tx:ex:fresh-small}
For $\Gamma=\{x,1-x\}$ over an arbitrary field, write a degree-one PC
proof of the fresh target $y$ and prove that its least NS degree is two.
Do not insert an axiom involving $y$.
\end{exercise}
\noindent{\small\hyperref[tx:ans:fresh-small]{Model answer (p.~\pageref*{tx:ans:fresh-small}).}}\par\medskip

\begin{exercise}\label{tx:ex:fresh-window}
In characteristic three, apply the fresh-$y$ bounds to the normalised
weak-bijection family at $n=6$ and $n=36$. What can the stated bounds
say at $n=9$? Distinguish the lower-bound size condition from the
all-size upper certificate.
\end{exercise}
\noindent{\small\hyperref[tx:ans:fresh-window]{Model answer (p.~\pageref*{tx:ans:fresh-window}).}}\par\medskip

\begin{exercise}\label{tx:ex:fresh-soundness}
Suppose $\Gamma$ has a root over some extension field of $K$. Explain
why the fresh target $y$ is underivable from $\Gamma^\uparrow$.
If one instead adds $y^2-y$ as an axiom, identify the extra term which
appears when extracting the coefficient of $y^1$ from an NS identity.
Why does the original coefficient-extraction proof no longer apply
unchanged?
\end{exercise}
\noindent{\small\hyperref[tx:ans:fresh-soundness]{Model answer (p.~\pageref*{tx:ans:fresh-soundness}).}}\par\medskip

